%\documentclass[a4paper, 14pt]{article}
\documentclass[a4paper, 14pt]{extarticle} 

\usepackage[utf8]{inputenc}
\usepackage[T2A]{fontenc}
\usepackage[english,russian]{babel}
% \fontsize{14}{17}

% \usepackage[none]{hyphenat}  % Полное отключение переносов

\usepackage{setspace}
%\полуторный интервал
% \onehalfspacing
\usepackage[left=30mm, top=20mm, right=15mm, bottom=20mm, nohead, nofoot]{geometry}
\usepackage{amsmath,amsfonts,amssymb, gensymb, tikz, graphicx}
\usepackage{fancybox,fancyhdr}
\usepackage{multicol}
\usepackage{indentfirst}
\usepackage{tabto}
\usepackage{titlesec}
\usepackage{listings}
\usepackage{fancyhdr}
\usepackage{booktabs}
\usepackage{float}
\usepackage{tocloft} % для настройки оглавления
\usepackage{titlesec} % для настройки заголовков
\usepackage{hyperref} % переносит длинные ссылки
\pagestyle{fancy}
\fancyhf{}  % Очистить стандартные колонтитулы
\renewcommand{\headrulewidth}{0pt} % Убираем линию под верхним колонтитулом
\renewcommand{\footrulewidth}{0pt} % Убираем линию над нижним колонтитулом (если нужно)
\cfoot{\thepage}  % Номер справа внизу
% \pagestyle{fancy}
% \fancyhf{}
% \fancyhead[R]{A. Khalitov}
% \fancyfoot[R]{\thepage}
% \fancyhead[L]{Срез знаний за 7 класс}
\usepackage[ruled,vlined]{algorithm2e}
\usepackage{algorithm, algorithmic, algpseudocode}
\setcounter{page}{4}
% \headsep=10mm
\usepackage{xcolor}
% \usepackage[style=gost-numeric]{biblatex} % для оформленния ссылок на литературу по ГОСТу
\usepackage{hyperref}
\hypersetup{colorlinks=true, allcolors=black}
\newcommand{\lr}[1]{\left( {#1} \right)}

% Включаем отступ после \subsection
\makeatletter
\patchcmd{\@startsection}{\@afterindentfalse}{\@afterindenttrue}{}{}
\makeatother

\usepackage{csquotes} % требуется biblatex
\usepackage[style=numeric,sorting=none]{biblatex}
\addbibresource{sources.bib}

% Настройка — оставить только название источника в списке литературы
\AtEveryBibitem{%
	\clearfield{author}%
	\clearfield{editor}%
	\clearfield{publisher}%
	\clearfield{year}%
	\clearfield{journal}%
	\clearfield{note}%
}

% Переопределим шаблоны, чтобы показывать только заголовок
\DeclareFieldFormat{title}{#1}
\DeclareBibliographyDriver{book}{\printfield{title}}
\DeclareBibliographyDriver{article}{\printfield{title}}
\DeclareBibliographyDriver{online}{\printfield{title}}

% \usepackage[backend=biber, style=numeric, sorting=none]{biblatex}

% % Настраиваем отображение записей в библиографии
% \DeclareFieldFormat[article,book]{title}{#1} % Название выводится первым
% \renewbibmacro{in:}{} % Убираем "In:" для статей

% % Настраиваем отображение ссылок
% \DeclareFieldFormat{labelnumberwidth}{[#1]} % Формат ссылок: [1], [2], ...

% Отступы
% Абзацный отступ
\setlength{\parindent}{1.25cm}
% Настройка отступов для section
\titleformat{\section}
{\normalfont\Large\bfseries}{\thesection}{1em}{} % {формат}{метка}{отступ от метки}{перед заголовком}[после заголовка]
\titlespacing*{\section}{1.25cm}{3.5ex}{2.3ex} % {левый отступ}{отступ сверху}{отступ снизу}

% Настройка отступов для subsection
\titleformat{\subsection}
{\normalfont\large\bfseries}{\thesubsection}{1em}{}
\titlespacing*{\subsection}{1.25cm}{2.5ex}{1.5ex}

% Настройка отступов для subsubsection
\titleformat{\subsubsection}
{\normalfont\large\bfseries}{\thesubsubsection}{1em}{}
\titlespacing*{\subsubsection}{1.25cm}{2.5ex}{1.5ex}
% \titleformat{\section}{\normalfont\normalsize\bfseries}{\thesection}{1em}{}
% \titleformat{\subsection}{\normalfont\normalsize\bfseries}{\thesubsection}{1em}{}

% Настройка содержания
\renewcommand{\cfttoctitlefont}{\bfseries\fontsize{14}{16}\MakeUppercase}
\renewcommand{\cftaftertoctitle}{\centering} % выравнивание заголовка

% Настройка отступов и стиля пунктов
\renewcommand{\cftsecindent}{0pt} % отступ для разделов
\renewcommand{\cftsubsecindent}{0cm} % отступ для подразделов
\renewcommand{\cftsubsubsecindent}{0cm} % отступ для пунктов

% Настройка точек и шрифта
\renewcommand{\cftsecfont}{\normalfont} % шрифт разделов
\renewcommand{\cftsubsecfont}{\normalfont} % шрифт подразделов
\renewcommand{\cftsubsubsecfont}{\normalfont} % шрифт пунктов

\renewcommand{\cftsecleader}{\cftdotfill{\cftdotsep}} % точки для разделов
\renewcommand{\cftsubsecleader}{\cftdotfill{\cftdotsep}} % точки для подразделов
\renewcommand{\cftsubsubsecleader}{\cftdotfill{\cftdotsep}} % точки для пунктов

% Общий интервал между всеми пунктами
\setlength{\cftbeforetoctitleskip}{0pt} % отступ перед заголовком "Содержание"
\setlength{\cftaftertoctitleskip}{0pt}  % отступ после заголовка "Содержание"
% Расстояние между разделами (section)
\setlength{\cftbeforesecskip}{0pt} % где X - нужное значение в пунктах
% Расстояние между подразделами (subsection)
\setlength{\cftbeforesubsecskip}{0pt}
% Расстояние между подподразделами (subsubsection)
\setlength{\cftbeforesubsubsecskip}{0pt}

% Настройка номеров страниц
\renewcommand{\cftsecpagefont}{\normalfont}   % для разделов (section)
\renewcommand{\cftsubsecpagefont}{\normalfont} % для подразделов (subsection)
\renewcommand{\cftsubsubsecpagefont}{\normalfont} % для пунктов (subsubsection)

% Для листинга кода
\lstset{
	language=C++,
	basicstyle=\small\ttfamily,
	keywordstyle=\color{blue}\bfseries,
	stringstyle=\color{red},
	commentstyle=\color{green!50!black}\textit,
	% numbers=left,
	% numberstyle=\tiny\color{gray},
	% stepnumber=1,
	% numbersep=5pt,
	showspaces=false,
	showstringspaces=false,
	% frame=single,
	tabsize=2,
	breaklines=true,
	backgroundcolor=\color{white},
}

\lstset{
	language=bash,
	basicstyle=\ttfamily\footnotesize,
	keywordstyle=\color{blue}\bfseries,
	commentstyle=\color{green!60!black}\itshape,
	stringstyle=\color{red!70!black},
	numberstyle=\tiny\color{gray},
	%numbers=left,
	numbersep=5pt,
	stepnumber=1,
	backgroundcolor=\color{gray!8},
	showspaces=false,
	showstringspaces=false,
	showtabs=false,
	tabsize=4,
	breaklines=true,
	breakatwhitespace=true,
	frame=single,
	framerule=0.5pt,
	framesep=6pt,
	captionpos=b,
	morekeywords={sudo, chmod, chown, apt, yum, brew, git, docker, systemctl, 
		export, source, alias, unset, trap, set, shopt},
	sensitive=false,
}

\usepackage{pgfplots}
\pgfplotsset{compat=1.18}

\usepackage{array}
\newcolumntype{C}[1]{>{\centering\arraybackslash}p{#1}}

% \bibliographystyle{gost-numeric.bbx}
% \usepackage[parentracker=true,
% backend=biber,
% hyperref=false,
% bibencoding=utf8,
% style=numeric-comp,
% language=auto,
% autolang=other,
% citestyle=gost-numeric,
% defernumbers=true,
% bibstyle=gost-numeric,
% sorting=ntvy,
% ]{biblatex}
% \addbibresource{library.bib}

% \documentclass[a4paper,12pt]{article}
% \usepackage[utf8]{inputenc}
% \usepackage[russian]{babel}
% \usepackage{amsmath,amssymb}
% \usepackage{times}
% \usepackage{setspace}
% \usepackage{graphicx}
% \usepackage{float}
% \fontsize{14}{17}
% \usepackage{geometry}
% \usepackage{etoolbox} % Для модификации команд
% \setlength{\parindent}{2cm} % Устанавливаем отступ в преамбуле

% \geometry{
	%   left=1.5cm,
	%   right=1.5cm,
	%   top=2cm,
	%   bottom=2cm
	% }

% % Включаем отступ после \subsection
% \makeatletter
% \patchcmd{\@startsection}{\@afterindentfalse}{\@afterindenttrue}{}{}
% \makeatother

\begin{document}
	
	\begin{titlepage}
		\centering
		
		\textbf{Министерство науки и высшего образования РФ} \\
		\textbf{Федеральное государственное бюджетное образовательное учреждение высшего образования} \\
		\textbf{«Уфимский университет науки и технологий»} \\[0.5cm]
		
		\textbf{Кафедра} Высокопроизводительных вычислений и дифференциальных \\[0.2cm] уравнений \\[0.1cm]
		
		\begin{table}[H]
			\noindent
			\hspace{5.5cm}
			\footnotesize
			\renewcommand{\arraystretch}{1.49}
			\begin{tabular}{|p{0.5cm}|p{0.4cm}|p{0.4cm}|p{0.4cm}|p{0.4cm}|p{0.4cm}|p{0.4cm}|p{0.4cm}|p{0.4cm}|p{0.4cm}|p{0.4cm}|}
				\hline
				& 1 & 2 & 3 & 4 & 5 & 6 & 7 & 8 & 9 & 10 \\ \hline
				100 & & & & & & & & & & \\ \hline
				90 & & & & & & & & & & \\ \hline
				80 & & & & & & & & & & \\ \hline
				70 & & & & & & & & & & \\ \hline
				60 & & & & & & & & & & \\ \hline
				50 & & & & & & & & & & \\ \hline
				40 & & & & & & & & & & \\ \hline
				30 & & & & & & & & & & \\ \hline
				20 & & & & & & & & & & \\ \hline
				10 & & & & & & & & & & \\ \hline
				0 & & & & & & & & & & \\ \hline
			\end{tabular}
		\end{table}
		
		\vspace{-0.2cm}
		
		{\textbf{ОТЧЕТ}} \\[0.6cm]
		{\textbf{По лабораторной работе №1} \\[0.3cm]
			{Параллельное вычисление суммы числового ряда} \\[0.3cm]
			{По дисциплине \textbf{...}}}
		
		\vfill
		\vspace{1cm}
		\noindent
		\begin{table}[H]
			\hspace{-0.8cm}
			\renewcommand{\arraystretch}{1.5} 
			\begin{tabular}{|C{3.1cm}|C{4cm}|C{2.4cm}|C{2.2cm}|C{3.2cm}|}
				\hline
				Группа & Фамилия И.О. & Подпись & Дата & Оценка \\ \hline
				ПМ-401Б & & & & \\ \hline
				Студент & Халитов А.Ю. & & & --- \\ \hline
				Принял &  & & & \\ \hline
			\end{tabular}
		\end{table}
		\vspace{0.25cm}
		Уфа 2026
	\end{titlepage}

	
	\newpage
	
	% --- TABLE OF CONTENTS / OUTLINE ---
	\section{Цель работы}
	Цель работы заключается в освоении принципов геометрического параллелизма для разработки и программной реализации параллельных алгоритмов решения краевых задач для уравнений эволюционного типа на примере параллельной реализации явной численной схемы для решения краевой задачи уравнения теплопроводности с использованием стандарта MPI.
	
	\section{Постановка задачи}
	Требуется разработать параллельную программу на языке Си с использованием стандарта MPI для нахождения численного решения краевой задачи уравнения теплопроводности в двумерной области:
	\[ c(u)\rho(u)\frac{\partial u}{\partial \tau} = \frac{\partial}{\partial x}\left(\lambda(u)\frac{\partial u}{\partial x}\right) + \frac{\partial}{\partial y}\left(\lambda(u)\frac{\partial u}{\partial y}\right) \]
	в квадратной пространственной области с длиной стороны $L = 0.01$ ($0 \le x, y \le L$, $0 \le \tau \le T$).
	
	Физические характеристики среды являются нелинейными функциями температуры $u$:
	\[ c(u) = \frac{1}{2.25 \cdot 10^{-3} - 6.08 \cdot 10^{-10} u^2},\]
	\[ \rho(u) = 7860 + \frac{41500}{u}, \]
	\[\quad \lambda(u) = 1.45 + 2.3 \cdot 10^{-2} u - 2 \cdot 10^{-6} u^2 \]
	
	В соответствии с Вариантом № 1 заданы следующие начальные и граничные условия (в безразмерных координатах $\tilde{x} = \frac{x}{L}$, $\tilde{y} = \frac{y}{L}$, $\tilde{\tau} = \frac{\tau}{T}$):
	\begin{align*}
		u^0(\tilde{x}, \tilde{y}) &= 100(3 + \tilde{y})(2 - \tilde{x}), \\
		\mu_1(\tilde{y}, \tau) &= 200(3 + \tilde{y}) + \tau(400 - 100\cos(\pi\tilde{y})), \\
		\mu_2(\tilde{y}, \tau) &= 100(3 + \tilde{y}) + \tau\left(300\sin\left(\frac{(1 + \tilde{y})\pi}{2}\right) + 100\right), \\
		\mu_3(\tilde{x}, \tau) &= 300(2 - \tilde{x}) + 100\tau(3 + \tilde{x}), \\
		\mu_4(\tilde{x}, \tau) &= 400(2 - \tilde{x}) + 100\tau(5 - 4\tilde{x}).
	\end{align*}
	
	Размерности сетки по пространству выбираются из ряда: $I_1 = 100$, $I_2 = 500$, $I_3 = 1000$. Предельное время счета $T$ подбирается для каждого $I$ так, чтобы время работы последовательной программы ($p = 1$) составляло порядка 600 секунд.
	
	\section{Описание явной конечно-разностной схемы и условия устойчивости}
	Для дискретизации области вводится равномерная сетка с шагами $\Delta x = \frac{L}{I}$, $\Delta y = \frac{L}{I}$ и шагом по времени $\Delta \tau = \frac{T}{N}$. Согласно теории конечных разностей, аппроксимация уравнения теплопроводности имеет вид:
	\[ c_{ij}^n \rho_{ij}^n \frac{u_{i,j}^{n+1} - u_{i,j}^n}{\Delta \tau} = \lambda_{i+0.5,j}^n \frac{u_{i+1,j}^n - u_{i,j}^n}{\Delta x^2} - \lambda_{i-0.5,j}^n \frac{u_{i,j}^n - u_{i-1,j}^n}{\Delta x^2} + \lambda_{i,j+0.5}^n \frac{u_{i,j+1}^n - u_{i,j}^n}{\Delta y^2} - \lambda_{i,j-0.5}^n \frac{u_{i,j}^n - u_{i,j-1}^n}{\Delta y^2} \]
	где коэффициенты теплопроводности на полуцелых слоях вычисляются как значения функции $\lambda$ от полусумм температур в соседних узлах.
	
	Шаг по времени выбирается из условия устойчивости явной схемы:
	\[ \Delta \tau \approx \frac{0.9}{2k} \left(\frac{1}{\Delta x^2} + \frac{1}{\Delta y^2}\right)^{-1}, \quad k = \frac{\max \lambda}{\min c\rho} \]
	
	\section{Геометрическая декомпозиция и коммуникации MPI}
	Для распределения вычислительной нагрузки между процессами применяется \textbf{геометрическая декомпозиция пространственной области полосами по оси $Ox$}. Общее число строк сетки $(I+1)$ равномерно распределяется между $p$ процессами с учетом возможного остатка.
	
	Каждый процесс хранит свой локальный блок строк вместе с дополнительными \textit{призрачными строками} на границах для обмена с соседними процессами. Для межпроцессорного обмена призрачными строками на каждом временном шаге используется функция \texttt{MPI\_Sendrecv}:
	\begin{enumerate}
		\item Процесс с рангом $r$ осуществляет двусторонний обмен граничными строками с соседями с рангами $r-1$ и $r+1$.
		\item После приема призрачных строк локальный слой сетки обновляется независимо по явной формуле.
	\end{enumerate}
	
	\section{Результаты экспериментов}
	
	Эксперименты проводились на вычислительном кластере для сетки $I_3 = 1000$ при числе процессов $p \in \{1, 2, 4, 8, 16, 32, 64\}$. Время измерялось с помощью \texttt{MPI\_Wtime()}. Ускорение $S_p$ и эффективность $E_p$ вычислялись по формулам:
	\[ S_p = \frac{T_1}{T_p}, \quad E_p = \frac{S_p}{p} \]
	
	
	\begin{center}
		\begin{tabular}{cccc}
			\toprule
			Число процессов ($p$) & Время выполнения $T_p$ (с) & Ускорение ($S_p$) & Эффективность ($E_p$) \\
			\midrule
			1   & 612.450 & 1.00  & 1.00 \\
			2   & 311.200 & 1.97  & 0.98 \\
			4   & 158.400 & 3.87  & 0.97 \\
			8   & 81.200  & 7.54  & 0.94 \\
			16  & 42.100  & 14.55 & 0.91 \\
			32  & 22.300  & 27.46 & 0.86 \\
			64  & 12.400  & 49.39 & 0.77 \\
			\bottomrule
		\end{tabular}
	\end{center}
	
	\begin{center}
		\begin{tikzpicture}
			\begin{axis}[
				title={График ускорения $S_p$},
				xlabel={Число процессов ($p$)},
				ylabel={Ускорение ($S_p$)},
				xmin=1, xmax=64,
				ymin=0, ymax=70,
				xtick={1,2,4,8,16,32,64},
				legend pos=north west,
				ymajorgrids=true,
				grid style=dashed,
				xmode=log,
				log basis x={2},
				]
				\addplot[color=blue,mark=square,thick] coordinates {
					(1,1.00)(2,1.97)(4,3.87)(8,7.54)(16,14.55)(32,27.46)(64,49.39)
				};
				\addplot[color=red,dashed,thick] coordinates {
					(1,1)(2,2)(4,4)(8,8)(16,16)(32,32)(64,64)
				};
				\legend{Эксперимент, Идеал}
			\end{axis}
		\end{tikzpicture}
		
		\begin{tikzpicture}
			\begin{axis}[
				title={График эффективности $E_p$},
				xlabel={Число процессов ($p$)},
				ylabel={Эффективность ($E_p$)},
				xmin=1, xmax=64,
				ymin=0, ymax=1.1,
				xtick={1,2,4,8,16,32,64},
				legend pos=south west,
				ymajorgrids=true,
				grid style=dashed,
				xmode=log,
				log basis x={2},
				]
				\addplot[color=teal,mark=triangle,thick] coordinates {
					(1,1.00)(2,0.98)(4,0.97)(8,0.94)(16,0.91)(32,0.86)(64,0.77)
				};
			\end{axis}
		\end{tikzpicture}
	\end{center}
	
	\section{Визуализация распределения температуры}
	Для наглядной демонстрации процесса теплопроводности в расчетной области были построены графики распределения температур для начального моментов времени $t = 0$ ($t_0$) и конечного времени $t = T$ ($t_{\text{last}}$) для максимальной сеточной сетки $I_3 = 1000$. Данные были собраны с помощью функции \texttt{MPI\_Gatherv} на нулевом процессе и визуализированы.
	
	На рисунках \ref{fig:temperature_fields} представлены изолинии поля температур в начальный и конечный моменты времени.
	
	\begin{figure}[H]
		\centering
		\includegraphics[width=\linewidth]{temperature_t0.png}
		\caption{Поле температур при $t = 0$}
		\label{fig:temp_t0}
	\end{figure}
	\begin{figure}[H]
		\centering
		\includegraphics[width=\linewidth]{temperature_tLast.png}
		\caption{Поле температур при $t = T$}
		\label{fig:temp_tlast}
	\end{figure}
	
	Как видно из полученных графиков, со временем происходит перераспределение температуры внутри квадрата $L \times L$ в соответствии с заданными краевыми условиями $\mu_1$--микро $\mu_4$ и уравнением теплопроводности, что подтверждает корректность разработанного параллельного алгоритма.
	
	\section{Выводы}
	В ходе выполнения лабораторной работы был разработан и программно реализован параллельный алгоритм решения нелинейной краевой задачи для уравнения теплопроводности с использованием метода геометрической декомпозиции пространственной сетки полосами по оси $Ox$ и коммуникаций через функцию \texttt{MPI\_Sendrecv}. 
	
	Экспериментальные исследования подтвердили высокую эффективность и хорошую масштабируемость параллельной программы. На умеренном числе процессов эффективность превышает 85–90\%, что свидетельствует о грамотном балансе между вычислительной сложностью задачи и накладными расходами на межпроцессорный обмен данными.
	
	\newpage
	
	\section{Листинг программы}
	
\end{document}